\documentclass[12pt]{article}
\input{structure.tex}
\renewcommand{\bottomtitlespace}{2cm}

\begin{document}

\renewcommand{\tl}{$3$ сек.}
\renewcommand{\ml}{$1024$ MB}
\problem{Задача A2. МУМИЯ}
\addauthor{Автор: Емил Инджев}{2025}{11}{22}
\renewcommand{\header}{
	НАЦИОНАЛЕН ЕСЕНЕН ТУРНИР ПО ИНФОРМАТИКА\\
	София, 21 – 23 ноември 2025 г.\\
	Група A – 11, 12 клас
}

Мумията Амуму играе срещу странен отбор от $N$ опонента. Всеки опонент е избрал да играе различен герой (номерирани с числата от $0$ до $N-1$) и ще играе различна роля (отново номерирани с  числата от $0$ до $N-1$). Амуму обаче няма идея кой герой в коя роля ще се играе, а би искал да знае. За целта той може да прави запитвания, в които дава свое пълно предположение за героя за всяка роля, и му се казва колко от тях е познал. 

Амуму обаче е просто една тъжна мумия, а не информатик. Затова Вие трябва да решите тази задача вместо него като напишете програма \textbf{\texttt{mummy}}. Формално, трябва да познаете скрита пермутация като правите заявки. Всяка заявка представлява пермутация, а като отговор получавате броя съвпадащи елементи със скритата пермутация. Целта е да минимизирате (колкото можете) средния брой заявки, които правите (все пак времето в champ select е лимитирано).

\subsection{Детайли по имплементацията}
Трябва да имплементирате функцията \texttt{findPerm}:
\begin{lstlisting}
   std::vector<int> findPerm(int N)
\end{lstlisting}
\begin{itemize}
	\item $N$: големина на пермутацията.
\end{itemize}

Тази функция ще бъде извикана $T$ пъти на тест -- по веднъж за всеки подтест, всички с еднакво $N$. Тя трябва да върне скритата пермутация. За целта програмата Ви може да прави заявки като извиква функцията \texttt{numMatches}:
\begin{lstlisting}
   int numMatches(const std::vector<int>& perm);
\end{lstlisting}
\begin{itemize}
	\item $perm$: пермутацията, за която е заявката.
\end{itemize}

Функцията приема валидна пермутация на числата от $0$ до $N-1$ и връща броя съвпадащи елементи със скритата пермутация.

\subsection{Ограничения}
\vspace{0.1em}
\begin{itemize}
\item $1 \leq N \leq 500$;
\item $T = 25$.
\end{itemize}

\subsection{Подзадачи и оценяване}
\begin{table}[H]
\begin{tblr}{|Q[c,m]|Q[c,m]|Q[c,m]|Q[c,m]|Q[c,m]|}
\hline
\textbf{Подзадача} & \textbf{Точки} & $N$ \\
\hline
$1$ & $9$ & $6$ \\
\hline
$2$ & $15$ & $44$ \\
\hline
$3$ & $7$ & $155$ \\
\hline
$4$ & $22$ & $300$ \\
\hline
$4$ & $47$ & $500$ \\
\hline
\end{tblr}
\end{table}
\FloatBarrier

Всяка подзадача се състои от точно един тест с $T=25$ подтеста с еднакво $N$. Точки за подзадача се получават само при валидно познаване на скритата пермутация за всеки подтест. Нека $Q$ е средния брой направени заявки в подзадачата. Тогава частта от точките за подзадачата, която ще получите е: $\frac{3700}{\max(Q, 3700)}$.

\subsection{Примерна интеракция}
В този пример $N = 4$, а скритата пермутация е $1, 3, 2, 0$.
\begin{table}[H]
\begin{tblr}{|Q[c,m]|Q[c,m]|Q[c,m]|Q[c,m]|Q[c,m]|}
\hline
\textbf{Вашата програма} & \textbf{Грейдъра} \\
\hline
 & \texttt{findPerm(4)} \\
\hline
\texttt{numMatches(\{0, 1, 2, 3\})} & \texttt{return 1} \\
\hline
\texttt{numMatches(\{3, 2, 1, 0\})} & \texttt{return 1} \\
\hline
\texttt{numMatches(\{0, 2, 3, 1\})} & \texttt{return 0} \\
\hline
\texttt{numMatches(\{1, 3, 0, 2\})} & \texttt{return 2} \\
\hline
\texttt{return \{1, 3, 2, 0\}} &  \\
\hline
\end{tblr}
\end{table}

Тук Вашата програма успява да познае скритата пермутация, използвайки 4 заявки.

\subsection{Локален грейдър}
Формат на входа:
\begin{itemize}
	\item ред $1$: две цели числа $N$ и $T$ -- големина на пермутациите и брой подтестове;
	\item ред $2$ до $T+1$: $N$ цели числа, описващи съответната пермутация.
\end{itemize}

Формат на изхода:
\begin{itemize}
	\item ред $1$: съобщение за грешка или средния брой заявки за подтестовете, ако всички пермутации са валидно познати.
\end{itemize}

\end{document}
